% Generated draft: requires mathematical and source review.
\documentclass[11pt]{article}
\usepackage[margin=25mm]{geometry}
\usepackage{fontspec}
\setmainfont{Latin Modern Roman}
\usepackage{amsmath,amssymb,mathtools,unicode-math}
\setmathfont{Latin Modern Math}
\usepackage{enumitem,needspace,xurl,hyperref}
\setlength{\parindent}{0pt}
\setlength{\parskip}{5pt}
\setlength{\emergencystretch}{3em}
\newcommand{\N}{\mathbb{N}}
\newcommand{\Z}{\mathbb{Z}}
\newcommand{\Q}{\mathbb{Q}}
\newcommand{\R}{\mathbb{R}}
\newcommand{\C}{\mathbb{C}}
\newcommand{\F}{\mathbb{F}}
\newcommand{\A}{\mathbb{A}}
\newcommand{\PP}{\mathbb{P}}
\newcommand{\E}{\mathbb{E}}
\newcommand{\eps}{\varepsilon}
\newcommand{\dd}{\,\mathrm{d}}
\begin{document}
\section*{UnsolvedMath ID: 30002018}
\textbf{Source problem code:} catalogue entry
\subsection{Definitions and mathematical statement}
The supplied statement concerns blocks with abelian defect groups for suitable finite groups of Lie type. Its intended scope and Deligne--Lusztig construction are not specified sufficiently to determine a fully quantified conjecture. The following gives a standard field-coefficient formulation of the terminology, without asserting that these choices exhaust the intended setting.

Let \(\ell\) be a prime, let \(k\) be an algebraically closed field of characteristic \(\ell\), and let \(G\) be a finite group of Lie type. A block idempotent \(b\) is a primitive central idempotent of \(kG\); its block algebra is \(A=kGb\). Suppose that \(D\leq G\) is an abelian defect group of this block. Let \(N_G(D)=\{g\in G:gDg^{-1}=D\}\), and let \(c\) be the Brauer correspondent of \(b\) in \(kN_G(D)\), with block algebra \(B=kN_G(D)c\). Write \(A\text{-}\mathrm{mod}\) and \(B\text{-}\mathrm{mod}\) for the categories of finite-dimensional left modules, and \(D^b\) for their bounded derived categories.

Here is an explicit convention for perversity. Let \(\mathcal S\) and \(\mathcal T\) be the sets of isomorphism classes of simple \(A\)-modules and simple \(B\)-modules. A triangulated equivalence
\[
F:D^b(A\text{-}\mathrm{mod})\longrightarrow D^b(B\text{-}\mathrm{mod})
\]
is perverse if there are finite filtrations
\[
\varnothing=\mathcal S_{-1}\subseteq\mathcal S_0\subseteq\cdots\subseteq\mathcal S_r=\mathcal S,
\qquad
\varnothing=\mathcal T_{-1}\subseteq\mathcal T_0\subseteq\cdots\subseteq\mathcal T_r=\mathcal T,
\]
and a function \(p:\{0,\ldots,r\}\to\Z\) with the following properties. For every \(i\) and every simple module \(S\) whose class belongs to \(\mathcal S_i\setminus\mathcal S_{i-1}\), all composition factors of \(H^j(F(S))\) belong to \(\mathcal T_{i-1}\) when \(j\ne p(i)\). All composition factors of \(H^{p(i)}(F(S))\) belong to \(\mathcal T_i\), and its image in the Serre quotient
\[
\frac{\{M\in B\text{-}\mathrm{mod}:\text{all composition factors of }M\text{ belong to }\mathcal T_i\}}
{\{M\in B\text{-}\mathrm{mod}:\text{all composition factors of }M\text{ belong to }\mathcal T_{i-1}\}}
\]
is simple. These simple images induce a bijection between \(\mathcal S_i\setminus\mathcal S_{i-1}\) and \(\mathcal T_i\setminus\mathcal T_{i-1}\). This convention places the distinguished cohomology in degree \(p(i)\).

\textbf{Requested assertion, subject to unresolved source qualifications.} For the suitable finite groups of Lie type and abelian-defect blocks intended in the source, show that the Deligne--Lusztig-derived equivalence predicted in connection with Brou\'e's abelian-defect-group conjecture is a perverse equivalence. In the field-coefficient formulation above, this asks for the intended functor \(F\) to satisfy the stated perversity conditions.

The supplied record does not identify the suitable groups, primes, blocks, coefficient setting, Deligne--Lusztig variety or complex, or intended filtrations and perversity function. It also does not specify whether existence of the predicted equivalence is assumed or must be established, or whether particular perversity data are prescribed. Consequently, this template must not be read as an unconditional assertion for every choice of \(G,\ell,k,b,D\).
\subsection{Short English statement}
The problem asks whether the Deligne--Lusztig equivalences predicted for suitable abelian-defect blocks of finite groups of Lie type are perverse: simple modules should correspond in filtered layers, with one distinguished cohomological degree in each layer. The supplied record leaves the eligible cases and the geometric construction unspecified, so a fully quantified statement requires further source clarification.
\subsection{Sources}
\begin{itemize}
\item Supplied local UnsolvedMath registry record, permanent ID \(30002018\), source problem code \texttt{catalogue entry}, titled \emph{Perverse Deligne--Lusztig Equivalences for Abelian Defect Blocks}, in the Oberwolfach Reports Open Problems collection. No individual proposer is identified; the supplied proposed year is 2012.
\item Supplied report citation: \emph{Representations of Finite Groups} (2012), Oberwolfach Reports, pp.~963--1019, DOI 10.4171/owr/2012/16. \url{https://doi.org/10.4171/owr/2012/16}. The record describes Oberwolfach Reports as documenting mathematical workshops and their open-problem discussions; it supplies no narrower page reference for this problem.
\item Supplied literature reference: \emph{Perverse equivalences and Brou\'e's conjecture}. \url{https://doi.org/10.1016/j.aim.2013.07.010}.
\item Supplied literature reference, labeled ``Cyclic case'': \url{https://arxiv.org/abs/1207.0116}.
\item Literature assessment checked at \texttt{2026-08-21}: the supplied curation labels the problem partially solved and states that perverse equivalences have been constructed in cyclic and several noncyclic cases, while the general equivalences predicted from Deligne--Lusztig varieties are not established. This is supplied dated curation derived from the OWR record, not an independently verified literature assessment; its status and cited claims require independent verification.
\end{itemize}
\end{document}
